\exercisegroup

\begin{enumerate}
\def\labelenumi{\arabic{enumi})}
\tightlist
\item
  \begin{enumerate}
  \def\labelenumii{\alph{enumii})}
  \tightlist
  \item
    设 \(\mathbf{E}\) 是一个 Hausdorff 局部凸空间，\(\mathbf{E}'\) 是它的对偶。证明下列性质等价：
  \end{enumerate}
\end{enumerate}

α) E 是次桶型的 (III, 第 44 页, exerc. 7)。

β) E' 中每个强有界子集都是等度连续的。

\(\gamma)\) \(E'\) 中每个强有界子集都是相对弱紧的，并且 \(E\) 的初始拓扑就是 \(\tau(E, E')\) 。

δ) 双对偶 \(E''\) 的强拓扑在 E 上诱导的拓扑与 E 上的初始拓扑相同。

于是，\(\mathbf{E}''\) 的强拓扑的 0 的一个基本邻域系，由对于拓扑 \(\sigma (\mathrm{E}'', \mathrm{E}')\) 而言的闭包组成，这些闭包来自 \(\mathbf{E}\) 的初始拓扑的 0 的一个基本邻域系。

\begin{enumerate}
\def\labelenumi{\alph{enumi})}
\setcounter{enumi}{1}
\tightlist
\item
  证明：若 E 是次桶型的，且它的对偶 \(E'\) 与它的代数对偶 \(E^{*}\) 相同，则 E 的初始拓扑是最细的局部凸拓扑。
\end{enumerate}

¶ 2) a) 证明次桶型空间的任何乘积都是次桶型的。（化归为 Hausdorff 次桶型空间的情形；然后利用 exerc. 1 以及 IV, 第 14 页的命题 15。）

\begin{enumerate}
\def\labelenumi{\alph{enumi})}
\setcounter{enumi}{1}
\tightlist
\item
  给出一个既非有界型亦非桶型的 Hausdorff 次桶型局部凸空间的例子。（按 III, 第 45 页, exerc. 16 的做法，将其中的桶型空间换成次桶型空间，并利用 a)。）
\end{enumerate}

\begin{enumerate}
\def\labelenumi{\arabic{enumi})}
\setcounter{enumi}{2}
\item
  设 E 是一个复 Hausdorff 局部凸空间，\(E_{0}\) 是 E 的底实局部凸空间，并设 \(E^{\prime}\) 与 \(E_{0}^{\prime}\) 分别是 E 与 \(E_{0}\) 的对偶。证明由 \(f \mapsto Rf\) 给出的、从 \(E^{\prime}\) 到 \(E_{0}^{\prime}\) 上的典范 R-线性映射对于 \(E^{\prime}\) 与 \(E_{0}^{\prime}\) 上的 S-拓扑是一个同胚，其中 S 是 E 的有界子集构成的任意一个集合。由此推出从双对偶 \(E^{\prime\prime}\) 到双对偶 \(E_{0}^{\prime\prime}\) 上的典范 R-线性映射的定义，该映射对于弱拓扑 \(\sigma(E^{\prime\prime}, E^{\prime})\) 与 \(\sigma(E_{0}^{\prime\prime}, E_{0}^{\prime})\) ，同样对于强拓扑 \(\beta(E^{\prime\prime}, E^{\prime})\) 与 \(\beta(E_{0}^{\prime\prime}, E_{0}^{\prime})\) ，都是一个同胚；在这个映射下，E（视为嵌入于 \(E^{\prime\prime}\) 中）变换为 \(E_{0}\)（视为嵌入于 \(E_{0}^{\prime\prime}\) 中）。
\item
  设 E 是一个 Hausdorff 次桶型局部凸空间。
\end{enumerate}

\begin{enumerate}
\def\labelenumi{\alph{enumi})}
\tightlist
\item
  证明：若 E 的强对偶 \(E_{b}^{\prime}\) 是有界型的，则 E 的完备化 \(\hat{E}\)（视为 \(E^{\prime*}\) 的一个向量子空间，III, 第 21 页, 定理 2）包含在 E 的双对偶 \(E^{\prime\prime}\) 中。
\item
  我们称 E 是区分的，如果 \(E^{\prime\prime}\) 中每个对于拓扑 \(\sigma(\mathrm{E}^{\prime\prime}, \mathrm{E})\) 有界的子集，都包含在 E 的某个有界子集（对于该拓扑）的闭包之中。证明：E 是区分的，必须且只需它的强对偶 \(E_{b}^{\prime}\) 是桶型的（参见 IV, 第 15 页, 命题 3）。
\end{enumerate}

\begin{enumerate}
\def\labelenumi{\arabic{enumi})}
\setcounter{enumi}{4}
\tightlist
\item
  设 \(\mathbf{E}\) 是一个 Hausdorff 局部凸空间，\(\mathbf{E}'\) 是它的对偶。
\end{enumerate}

\begin{enumerate}
\def\labelenumi{\alph{enumi})}
\item
  要使 \(\mathbf{E}'\) 上的强拓扑与 \(\tau (\mathrm{E}',\mathrm{E})\) 相同，必须且只需 \(\mathbf{E}\) 是半自反的。
\item
  假设 E 是次桶型的。要使 \(E'\) 上的强拓扑与紧收敛拓扑相同，必须且只需 E 是一个 Montel 空间。
\end{enumerate}

¶ 6) 设 F 与 G 是处于分离对偶中的两个向量空间。

\begin{enumerate}
\def\labelenumi{\alph{enumi})}
\tightlist
\item
  证明下列性质等价：
\end{enumerate}

\(\alpha)\) 空间 \(\mathbf{F}\) ，带上一个与 \(\mathbf{F}\) 和 \(\mathbf{G}\) 之间的对偶相容的拓扑，是半自反的；

\(\beta)\) 空间 \(\mathbf{G}\) ，带上 \(\tau (\mathbf{G},\mathbf{F})\) ，是桶型的。

\begin{enumerate}
\def\labelenumi{\alph{enumi})}
\setcounter{enumi}{1}
\tightlist
\item
  证明下列性质等价：
\end{enumerate}

α) 空间 F，带上 τ(F, G)，是自反的；

\(\beta)\) 空间 \(\mathbf{G}\) ，带上 \(\tau (\mathbf{G},\mathbf{F})\) ，是自反的；

γ) F 与 G 分别对于 τ(F, G) 与 τ(G, F) 是桶型的。

\begin{enumerate}
\def\labelenumi{\alph{enumi})}
\setcounter{enumi}{2}
\item
  证明：每个 Hausdorff 局部凸空间 E，带上拓扑 \(\tau(\mathrm{E}, \mathrm{E}')\) ，同构于一个半自反空间 F 关于一个闭子空间的商。（由 III, 第 44 页, exerc. 14, c)，\(E'\) 带上 \(\tau(\mathrm{E}', \mathrm{E})\) 后同构于某个 Hausdorff 桶型空间 G 的一个闭向量子空间 M；取 \(F = G'\) 带上 \(\tau(\mathrm{G}', \mathrm{G})\) ，并利用 IV, 第 10 页的命题 11。）
\item
  设 A 是一个无限集，满足 \(\operatorname{Card}(\mathrm{A}) > \aleph_{1}\) (S, III, §6, exerc. 10)。在乘积空间 \(P = R^{A}\) 中，设 \(E_{0}\) 表示由所有满足如下条件的 \(x = (x_{\alpha})_{\alpha \in \mathrm{A}}\) 组成的处处稠密子空间：除对一个可数的指标集外，均有 \(x_{\alpha} = 0\) ；空间 \(E_{0}\) 是桶型的，并且 P 中由 \(E_{0}\) 与点 \(1_{A}\) （P 中坐标全等于 1 的点）生成的子空间 E 也是桶型的 (III, 第 45 页, exerc. 16)。设 B 是 \(E_{0}\) 中的一个有界子集，它在 P 中的闭包包含 \(1_{A}\) ，并设 J 是 A 的一个基数为 \(\aleph_{1}\) 的子集；设 \(pr_{J}\) 表示从 P 到 \(R^{J}\) 上的投影；证明：B 中存在一个子集 \(B_{J}\) ，其基数 \(\leqslant \aleph_{1}\) ，使得 \(\operatorname{pr}_{J}(B_{J})\) 在 \(R^{J}\) 中的闭包包含 \(1_{J}\) ；于是，集 \(J' \supset J\) 由所有满足如下条件的指标 \(\alpha \in A\) 组成：\(B_{J}\) 中至少有一点，其下标 \(\alpha\) 的坐标 \(\neq 0\) ；该集的基数等于 \(\aleph_{1}\) 。我们定义 \(J_{0} = J\) ，并用归纳法定义 \(J_{n+1} = J'_{n}\) ，令 \(H = \bigcup_{n} J_{n}\) （其基数为 \(\aleph_{1}\) ），并令 \(B_{H} = \bigcup_{n} B_{J_{n}}\) ；证明：\(B_{H}\) 的闭包（从而 B 的闭包）包含点 \((1_{H}, 0) \in R^{H} \times R^{A-H} = P\) ，而该点不属于 \(E_{0}\) 。由此推断：对 E 中每个有界紧集 C，\(C \cap E_{0}\) 在 E 中仍是闭的，并且 E 不是半自反的。
\item
  由 \(d\) 推出：\(\mathrm{E}'\) （即 \(\mathrm{E}\) 的对偶，它可与对偶 \(\mathrm{P}' = \mathbf{R}^{(\mathbf{A})}\) （\(\mathrm{P}\) 的对偶）等同）对于拓扑 \(\tau (\mathrm{E}',\mathrm{E})\) 不是完备的，尽管它对于该拓扑是半自反的（从而是拟完备的）（考虑 \(\mathrm{E}\) 上在 \(\mathrm{E}_0\) 上等于 0 而在点 \(1_{\mathrm{A}}\) 处等于 1 的线性型，并利用 III, 第 21 页, 定理 2）。
\end{enumerate}

\begin{enumerate}
\def\labelenumi{\arabic{enumi})}
\setcounter{enumi}{6}
\item
  设 \((\mathrm{E}_{i})_{i\in\mathbb{I}}\) 是一族 Hausdorff 局部凸空间，P 是诸 \(E_{i}\) 的乘积空间，S 是它们的拓扑直和。证明：要使 P 或 S 是半自反的（相应地，自反的），必须且只需每个 \(E_{i}\) 都是半自反的（相应地，自反的）。
\item
  设 E 是一个 Hausdorff 局部凸空间，它是一个闭向量子空间递增序列 \((E_{n})\) 的严格归纳极限 (II, 第 32 页, 命题 9)。
\end{enumerate}

\begin{enumerate}
\def\labelenumi{\alph{enumi})}
\item
  证明：若每个 \(\mathbf{E}_n\) 的强对偶都是完备的，则 \(\mathbf{E}\) 的强对偶是完备的 (III, 第 20 页, 定理 1)。
\item
  要使 E 是半自反的（相应地，自反的），必须且只需每个 \(E_{n}\) 都是半自反的（相应地，自反的）。
\end{enumerate}

\begin{enumerate}
\def\labelenumi{\arabic{enumi})}
\setcounter{enumi}{8}
\item
  设 \((\mathrm{E}_{\alpha})_{\alpha\in\mathrm{A}}\) 是包含于同一向量空间中的一族 Hausdorff 局部凸空间，它对于关系 \(\supset\) 是有向的，并且当 \(E_{\beta} \subset E_{\alpha}\) 时，\(E_{\beta}\) 的拓扑比 \(E_{\beta}\) 上由 \(E_{\alpha}\) 的拓扑诱导的拓扑更细。设 E 是诸 \(E_{\alpha}\) 的交，带上一个由诸 \(E_{\alpha}\) 的拓扑在 E 上诱导的拓扑取上确界所得的拓扑。证明：若每个 \(E_{\alpha}\) 都是半自反的，则 E 是半自反的（考虑 E 的一个有界子集上的一个超滤子）。
\item
  证明：Montel 空间的任何乘积以及任何拓扑直和都是一个 Montel 空间 \(^{1}\) 。
\item
  设 E 是一个 Hausdorff 局部凸空间，且 E 的强对偶 \(E_{b}^{\prime}\) 是半自反的。a) 证明：在 \(E^{\prime}\) 的每个强有界子集上，由 \(\sigma(E^{\prime}, E^{\prime\prime})\) 与 \(\sigma(E^{\prime}, E)\) 诱导的拓扑相同。
\end{enumerate}

\begin{enumerate}
\def\labelenumi{\alph{enumi})}
\setcounter{enumi}{1}
\tightlist
\item
  由 a) 推出：E 对于拓扑 \(\tau(E, E')\) 是次桶型的（参见 IV, 第 52 页, exerc. 1），并且若 \(\hat{E}\) 是 E 对于该拓扑的完备化，且视为 \(E'^{*}\) 的一个子集 (III, 第 21 页, 定理 2)，则 \(E'' \subset \hat{E}\) 。特别地，若 E 对于 \(\tau(E, E')\) 是拟完备的，则 E 对于该拓扑是自反的。
\end{enumerate}

\(^{1}\) 另一方面，Montel 空间的闭子空间未必是次桶型的，并且 Montel 空间关于闭子空间的商未必是半自反的 (IV, 第 63 页, exerc. 8)。

\begin{enumerate}
\def\labelenumi{\arabic{enumi})}
\setcounter{enumi}{11}
\tightlist
\item
  设 \(\mathbf{E}\) 是一个 Banach 空间，\(\mathbf{E}'\) 是它的对偶。
\end{enumerate}

\begin{enumerate}
\def\labelenumi{\alph{enumi})}
\item
  证明：由 \(x \mapsto d(x, \mathbf{A})\) 给出的、一点 \(x \in \mathbf{E}\) 到一个闭凸集 \(\mathbf{A}\) 的距离，是 \(\mathbf{E}\) 上对于拓扑 \(\sigma(\mathbf{E}, \mathbf{E}')\) 的下半连续函数。
\item
  证明：若 \(\mathrm{E}\) 是自反的，则对每个闭凸子集 \(\mathrm{A}\) （含于 \(\mathrm{E}\) ），存在一点 \(x_0 \in \mathrm{A}\) 使得 \(\| x_0 \|\) 等于 0 到 \(\mathrm{A}\) 的距离（利用 \(a\) ）。若 \(\mathrm{E}\) 的单位球的每个边界点都是极点，则该点是唯一的 (II, 第 54 页, 定义 1)。
\item
  假设 E 是自反的，并设 B 是 E 中一个闭的、凸的且有界的子集；由 a) 与 b) 推出：存在两点 \(x \in \mathbf{A}\) ，\(y \in \mathbf{B}\) ，使得 \(\| x - y \| = d(\mathbf{A}, \mathbf{B})\) （参见 V, 第 71 页, exerc. 8）。
\end{enumerate}

\begin{enumerate}
\def\labelenumi{\arabic{enumi})}
\setcounter{enumi}{12}
\item
  设 \(\mathbf{E}\) 是一个 Banach 空间，\(\mathbf{M}\) 是 \(\mathbf{E}\) 的一个闭向量子空间。证明：若 \(\mathbf{M}\) 与 \(\mathrm{E / M}\) 都是自反的，则 \(\mathbf{E}\) 是自反的。
\item
  设 \(\mathbf{A}\) 是一个无限集。
\end{enumerate}

\begin{enumerate}
\def\labelenumi{\alph{enumi})}
\tightlist
\item
  证明：Banach 空间 \(E = \overline{\mathcal{K}(A)}\) 的强对偶 (IV, 第 47 页, exerc. 1) 可以与 Banach 空间 \(\ell^{1}(A)\) 等同，并且 Banach 空间 \(\ell'(A)\) 的强对偶可以与 Banach 空间 \(\mathcal{B}(A) = \ell^{\infty}(A)\) 等同；由此推出 E 不是自反的，并且 \(E''/E\) 是无限维的（参见 IV, 第 71 页, exerc. 18）。若 A = N，则 E 与 \(E'\) 都是满足第一可数公理的 Banach 空间，但 \(E''\) 不是 (I, 第 25 页, exerc. 1)。
\item
  设 \(B''\) 是 \(E'' = \ell^{\infty}(A)\) 中的单位球，并设 \(B_{0}''\) 是凸集 \(B'' + (B'' \cap E)\) 。证明：\(B_{0}''\) 是 \(E''\) 中对于强拓扑的一个有界闭凸集，其内部非空，但它没有任何极点。若 p 是 \(B_{0}''\) 的度规，证明：带上范数 p 的 \(E''\) 不等距于任何 Banach 空间的对偶，并且 \(B_{0}''\) 对于拓扑 \(\sigma(E'', E')\) 不是闭的（尽管 \(B''\) 对于 \(\sigma(E'', E')\) 是紧的，且 \(B'' \cap E\) 是强闭的）。
\end{enumerate}

¶ 15) 记号与 exerc. 14 中相同。

\begin{enumerate}
\def\labelenumi{\alph{enumi})}
\tightlist
\item
  设 \((x_n')\) 是 \(E'\) 中对于拓扑 \(\sigma(E', E'')\) 收敛于 0 的一个序列；证明：对每个 \(\varepsilon > 0\) ，存在一个有限子集 \(H\) （\(A\) 的），使得 \(\sum_{\alpha \notin H} |x_n'(\alpha)| \leqslant \varepsilon\) 对每个整数 \(n\) 都成立。
\end{enumerate}

（用反证法：若该性质不成立，证明此时存在一个数 \(\delta > 0\) 、整数的一个递增序列 \((n_k)\) 、A 的有限子集的一个递增序列 \((\mathrm{H}_k)\) ，使得 \(\sum_{\alpha \in \mathrm{H}_{k-1}} |x_n'(\alpha)| \leqslant \frac{\delta}{8}\) 对 \(n \geqslant n_k\) 成立，\(\sum_{\alpha \notin \mathrm{H}_k} |x_n'(\alpha)| \leqslant \frac{\delta}{8}\) 对 \(n \leqslant n_k\) 成立，并且

\(\sum_{\alpha \in \mathrm{H}_k - \mathrm{H}_{k-1}} |x_n'(\alpha)| \geqslant \frac{\delta}{2}\) ；证明这导致矛盾（“滑动峰”方法）。）

由此推出序列 \((x_{n}^{\prime})\) 对于强拓扑收敛于 0，尽管强拓扑严格细于拓扑 \(\sigma (\mathrm{E}',\mathrm{E}'')\) 。

\begin{enumerate}
\def\labelenumi{\alph{enumi})}
\setcounter{enumi}{1}
\tightlist
\item
  证明：若 \((x_{n}^{\prime})\) 是 \(E'\) 中对于拓扑 \(\sigma(E', E'')\) 的一个 Cauchy 序列，则它对于该拓扑收敛于 \(E'\) 中的一点；换言之，\(E'\) 对于 \(\sigma(E', E'')\) 是半完备的 (III, 第 7 页) 。（证明：对每个 \(\varepsilon > 0\) ，存在一个有限子集 \(H\) （\(A\) 中的），使得对每个整数 n 都有 \(\sum_{\alpha \notin H} |x_{n}'(\alpha)| \leqslant \varepsilon\) ；按 a) 的方式用反证法，并利用 a)。）
\end{enumerate}

¶ 16) 采用 exerc. 14 的记号，设 E''' 是 \(E'' = \ell^{\infty}(A)\) 的对偶。

\begin{enumerate}
\def\labelenumi{\alph{enumi})}
\item
  设 \(e_{\alpha}\) （\(\alpha \in A\) ）是 \(E''\) 中满足 \(e_{\alpha}(\beta) = \delta_{\alpha \beta}\) （Kronecker 符号）的元素。设 \((\mathbf{K}_n)\) 是 \(A\) 的有限子集的一个序列，两两不相交，并设 \((x_n''')\) 是 \(E'''\) 中的点的一个序列。证明：存在一个严格递增的无穷整数序列 \((n_k)\) （\(>0\) ），使得若令 \(B = \bigcup_k K_{n_k}\) ，则元素 \(y_n' = (x_n'''(\alpha))_{\alpha \in B}\) 属于 \(\ell^1(B)\) 。（设 \(\delta\) 是任意一个 \(>0\) 的数，并设 \((\mathrm{J}_m)_{m \in \mathbb{N}}\) 是 \(\mathbf{N}\) 分成有限集的一个分割。用反证法证明：若 \(x''' \in E'''\) ，则存在一个整数 \(m\) ，使得 \(|\langle x'', x''' \rangle| \leqslant \delta\) 对所有满足下列条件的 \(x'' \in E''\) 成立：\(\| x'' \| \leqslant 1\) 且 \(x''(\alpha) = 0\) ，只要指标 \(\alpha\) 不属于集 \(\bigcup_{n \in J_m} K_n\) 。以适当的方式相继对 \(x_1'', x_2'', ...\) 应用这一结果。）
\item
  由 \(a\) 以及 exerc. 15, a) 推出：若 \((x_{n}^{\prime \prime \prime})\) 是 \(\mathrm{E}'''\) 中对于拓扑 \(\sigma (\mathrm{E}''',\mathrm{E}'')\) 收敛于 0 的一个序列，且 \(\tilde{x}_n^{\prime \prime \prime}\) 是 \(x_{n}^{\prime \prime \prime}\) 在 \(\mathrm{E}''\) 的强闭子空间 E 上的限制，则 \(\lim_{n\to \infty}\| \tilde{x}_n^{\prime \prime \prime}\| = 0\) （在 \(\mathrm{E}'\) 中）。
\item
  由 b) 推出：\(E''\) 的强闭子空间 E 对于强拓扑没有拓扑补。（限于 A = N 的情形，设 \((e_{n}')\) 是 E 上的连续线性型的一个序列，满足 \(\langle x, e_{n}' \rangle = x(n)\) 对所有 \(x \in E\) ；证明序列 \((e_{n}')\) 对于 \(\sigma(\mathrm{E}', \mathrm{E})\) 趋于 0，但 \(e_{n}'\) 不能扩张为连续线性型 \(x_{n}'''\) （\(E''\) 上的），使得序列 \((x_{n}'''')\) 对于 \(\sigma(\mathrm{E}'''', \mathrm{E}'')\) 趋于 0。）
\end{enumerate}

\begin{enumerate}
\def\labelenumi{\arabic{enumi})}
\setcounter{enumi}{16}
\tightlist
\item
  设 E 是一个非自反 Banach 空间，\(E'\) 是它的强对偶，\(E''\) 是 \(E'\) 的强对偶，\(E'''\) 是 \(E''\) 的强对偶，\(E^{IV}\) 是 \(E'''\) 的强对偶。
\end{enumerate}

\begin{enumerate}
\def\labelenumi{\alph{enumi})}
\item
  证明：在 \(E^{\prime\prime}\) 中，\(E^{\prime}\) 与 E 的正交子空间 \(E^{\circ}\) （当 E 视为 \(E^{\prime\prime}\) 的子空间时）是拓扑补，并且 \(E^{\prime\prime}\) 到 \(E^{\prime}\) 上（对于这个分解）的投影是范数为 1 的连续线性映射。与 exerc. 16, c) 相比较，推出 Banach 空间 \(\mathcal{K}(A)\) （作为拓扑向量空间）不同构于任何 Banach 空间的强对偶。
\item
  证明：\(E^{IV}\) 是 \(E^{\circ}\) 与 \(E''\) 的拓扑直和，也是 \(E^{\prime\circ}\) 与 \(E^{\circ\circ}\) 的拓扑直和；我们有 \(E'' \cap E^{\circ\circ} = E\) 。设 v 是 \(E^{IV}\) 到其自身上的线性映射，它在 \(E^{\circ}\) 上是恒同，而在 \(E''\) 上是 \(E''\) 到 \(E^{\circ\circ}\) 上平行于 \(E^{\circ}\) 的投影；证明 v 是一个等距，但对于拓扑 \(\sigma(\mathrm{E}^{\mathrm{IV}}, \mathrm{E}''')\) 不连续。
\end{enumerate}

\begin{enumerate}
\def\labelenumi{\arabic{enumi})}
\setcounter{enumi}{17}
\item
  证明：一个 Banach 空间 E，若其强对偶 \(E'\) 满足第一可数公理，且 E 对于弱化拓扑 \(\sigma(\mathrm{E}, \mathrm{E}')\) 是半完备的 (III, 第 7 页) ，则 E 是自反的（与 IV, 第 54 页, exerc. 15, b) 相比较）。
\item
  设 \(\mathrm{E}\) 是一个 Banach 空间，\(\mathrm{E}'\) 是它的强对偶，\(\mathrm{G}'\) 是 \(\mathrm{E}'\) 的一个对于强拓扑满足第一可数公理的强闭子空间。证明：存在 \(\mathrm{E}\) 的一个可数子集，使得若 \(\mathrm{F}\) 是该子集在 \(\mathrm{E}\) 中生成的闭向量子空间，则 \(\mathrm{G}'\) 等距于 \(\mathrm{F}'\) （\(\mathrm{F}\) 的强对偶）的一个强闭子空间。（设序列 \((x_n')\) 在 \(\mathrm{G}'\) 中强稠密；对每个 \(n\) ，取 \(x_n \in \mathrm{E}\) 使得 \(\| x_n \| \leqslant 1\) 且 \(\langle x_n, x_n' \rangle = \left(1 - \frac{1}{n}\right) \| x_n' \|\) ；证明：强闭子空间 \(F\) 在 \(E\) 中由诸
\end{enumerate}

\[
x _ {n}
\]

生成的强闭子空间 F 即为所求的空间。）

¶ 20) 设 E 是一个 Banach 空间，E' 是它的对偶，B 是 E 中的单位球。要使 B 的每一点对于 B 上由弱化拓扑 σ(E, E') 诱导的拓扑都具有可数的基本邻域系，必须且只需 E' 对于强拓扑满足第一可数公理。（为证明条件是必要的，注意：若 B 的每一点对于弱化拓扑都有可数的基本邻域系，则 B 在 E'' 中对于拓扑 σ(E'', E') 的闭包 B°° 的每一点也是如此。于是存在 E' 中的一个序列 (a'\_n) ，使得 B°° 中 0 的每个 σ(E'', E') 邻域都包含 B°° 与有限多个极集 \{a'\_n\}° 的交；考虑由诸 a'\_n 生成的 E' 的强闭子空间 W' ，以及 W' 在 E'' 中的正交 \(W'^{\circ}\) 。）

\begin{enumerate}
\def\labelenumi{\arabic{enumi})}
\setcounter{enumi}{20}
\tightlist
\item
  设 E 是一个 Banach 空间，\(E'\) 是它的对偶，B 是 E 中的单位球，\(B_{r}'\) 是 \(E'\) 中以 0 为中心、半径为 r 的闭球。
\end{enumerate}

\begin{enumerate}
\def\labelenumi{\alph{enumi})}
\item
  设 \(M_{1}^{\prime}\) ，\(M_{2}^{\prime}\) 是 \(E^{\prime}\) 的两个向量子空间，它们对于弱拓扑 \(\sigma(\mathrm{E}^{\prime}, \mathrm{E})\) 都是处处稠密的。要使 \(\sigma(\mathrm{E}, \mathrm{M}_{1}^{\prime})\) 与 \(\sigma(\mathrm{E}, \mathrm{M}_{2}^{\prime})\) 在 B 上诱导的拓扑相同，必须且只需 \(M_{1}^{\prime}\) 与 \(M_{2}^{\prime}\) 在 \(E^{\prime}\) 中的强闭包相同。
\item
  设 \(M'\) 是 \(E'\) 的一个对于 \(\sigma(E', E)\) 处处稠密的向量子空间；用 \(\mathbf{M}'^{(1)}\) 表示 \(M' \cap B_{1}'\) 在 \(E'\) 中对于拓扑 \(\sigma(E', E)\) 的闭包所生成的向量子空间。要使 \(\mathbf{M}'^{(1)} = E'\) ，必须且只需 \(M' \cap B_{1}'\) 的弱闭包包含一个球 \(B_{r}'\) ，其中 r \textgreater{} 0（利用 \(E'\) 对于强拓扑是桶型的这一事实）。
\item
  设 \(r\) 是这样的数 \(t\) 的上确界：\(\mathbf{M}' \cap \mathbf{B}_1'\) 的弱闭包包含一个球 \(\mathbf{B}_t'\) ；称数 \(r\) 为 \(\mathbf{M}'\) 的特征数。证明：\(r\) 是诸数 \(\sup_{x' \in \mathbf{M}' \cap \mathbf{B}_1'} |\langle x, x' \rangle| / \| x\|\) 的下确界，其中 \(x\) 取遍所有 \(\neq 0\) 且含于 \(E\) 的点（利用 Hahn-Banach 定理）。
\item
  证明：\(1 / r\) 是 \(\| x\|\) 的上确界，其中 \(x\) 取遍 B 在 E 中对于拓扑 \(\sigma (\mathrm{E},\mathbf{M}^{\prime})\) 的闭包（利用 \(c\) ) 与 Hahn-Banach 定理）。
\item
  设 \(M^{\prime\circ}\) 是 \(M^{\prime}\) 在 \(E^{\prime\prime}\) 中的正交；证明：\(r = \inf(\|x + z^{\prime\prime}\|/\|x\|)\) ，其中 \(z^{\prime\prime}\) 取遍 \(M^{\prime\circ}\) ，x 取遍 E 中非零的点（利用 c) 与 Hahn-Banach 定理）。由此推出：要使 \(\mathbf{M}^{\prime(1)} = \mathbf{E}^{\prime}\) ，必须且只需 \(E + M^{\prime\circ}\) 在 \(E^{\prime\prime}\) 中是强闭的（利用 I, 第 17 页的 Banach 定理）。
\end{enumerate}

\(f\) ) 设 \(\mathbf{A} = \mathbf{N}\times \mathbf{N}\) ，\(\mathrm{E} = \overline{\mathcal{K}(\mathbf{A})}\) （参见 IV, 第 54 页, exerc. 14）。在空间 \(\mathrm{E}'' = \ell^{\infty}(\mathrm{A})\) 中，设 \(\mathrm{P}\) 是由所有这样的点 \(x = (x_{ij})\) 组成的向量子空间：\(x_{ij} = x_{0j} / (j + 1)\) 对所有 \(i\geqslant 0\) 。证明：\(\mathbf{P} = \mathbf{M}'^{\circ}\) ，其中 \(\mathbf{M}'\) 是 \(\mathrm{E}'\) 的一个（对于 \(\sigma (\mathrm{E}',\mathrm{E}))\) 处处稠密的向量子空间，但 \(\mathrm{E} + \mathrm{M}'^{\circ} = \mathrm{E} + \mathrm{P}\) 在 \(\mathrm{E}''\) 中不是强闭的；由此推出 \(\mathbf{M}'\) 的特征数为 0。

\begin{enumerate}
\def\labelenumi{\arabic{enumi})}
\setcounter{enumi}{21}
\tightlist
\item
  设 E 是一个 Banach 空间，\(E'\) 是它的对偶，\(M'\) 是 \(E'\) 中的一个强闭子空间，它对于 \(E'\) 上的弱拓扑是处处稠密的。我们称 M' 是不可约的，如果不存在向量子空间 \(N' \neq M'\) （\(M'\) 的），它在 \(E'\) 中是强闭且弱处处稠密的。a) 证明：\(M'\) 不可约的充要条件是正交 \(M'^{\circ}\) （\(M'\) 在 \(E''\) 中的）是 E 的拓扑补（对于 \(E''\) 的强拓扑）。由此推出此时 \(\mathbf{M}^{(1)} = \mathbf{E}'\) (exerc. 21)，并且 E 同构于空间 \(M'\) （带上由 \(E'\) 的强拓扑诱导的拓扑）的强对偶。
\end{enumerate}

\begin{enumerate}
\def\labelenumi{\alph{enumi})}
\setcounter{enumi}{1}
\item
  证明：\(\mathbf{M}'\) 不可约的充要条件是 \(\mathbf{E}\) 中的单位球对于拓扑 \(\sigma (\mathrm{E},\mathbf{M}')\) 是相对紧的（利用 exerc. 21, a)）。
\item
  要使 E （作为拓扑向量空间）同构于某个 Banach 空间的强对偶，必须且只需 \(E'\) 中存在一个不可约子空间。由此给出下述事实的一个新证明：Banach 空间 \(\overline{\mathcal{K}(\mathbf{N})}\) 不同构于任何 Banach 空间的强对偶（参见 IV, 第 55 页, exerc. 16, c)）。
\end{enumerate}

\begin{enumerate}
\def\labelenumi{\arabic{enumi})}
\setcounter{enumi}{22}
\tightlist
\item
  采用与 exerc. 22 相同的记号，假设 \(\mathbf{M}'\) 是不可约的。
\end{enumerate}

\begin{enumerate}
\def\labelenumi{\alph{enumi})}
\item
  要使从 E 到 \(E^{\prime\prime}/M'^{\circ}\) 中的典范映射（它是典范映射 \(E^{\prime\prime} \rightarrow E^{\prime\prime}/M'^{\circ}\) 的限制）是一个 Banach 空间等距，必须且只需 \(M^{\prime}\) 的特征数 (IV, 第 55 页, exerc. 21) 等于 1。这时，称 \(M^{\prime}\) （带上由 \(E^{\prime}\) 的范数诱导的范数）为 E 的预对偶，并且 E 可以（连同其范数）典范地等同于 Banach 空间 \(M^{\prime}\) 的对偶。
\item
  对 E 的每个对于 \(\sigma (\mathrm{E},\mathbf{M}^{\prime})\) 为闭的向量子空间 F ，证明：\(\mathbf{M}'\) 在 F 的对偶 \(\mathrm{F}'\) （它等同于 \(\mathrm{E}^{\prime} / \mathrm{F}^{\circ}\) ）中的典范像是 F 的一个预对偶。
\end{enumerate}

\begin{enumerate}
\def\labelenumi{\arabic{enumi})}
\setcounter{enumi}{23}
\tightlist
\item
  \begin{enumerate}
  \def\labelenumii{\alph{enumii})}
  \tightlist
  \item
    设 \((a_{k})_{1\leqslant k\leqslant n}\) 是赋范空间 \(E\) 中点的一个有限序列，并设 \((\lambda_k)_{1\leqslant k\leqslant n}\) 是数的一个有限序列，诸数 \(>0\) 且满足 \(\sum_{k = 1}^{n}\lambda_{k} < 1\) ；令 \(\mu_{k} = 1 - \sum_{j = 1}^{k - 1}\lambda_{j}\) （对每个 \(k\) ）；那么
  \end{enumerate}
\end{enumerate}

\[
\left\| \sum_ {k = 1} ^ {n} \lambda_ {k} a _ {k} \right\| \leqslant \frac {\lambda_ {n}}{\mu_ {n - 1}} \left\| \mu_ {n - 1} a _ {n} + \sum_ {k = 1} ^ {n - 1} \lambda_ {k} a _ {k} \right\| + \frac {\mu_ {n}}{\mu_ {n - 1}} \left\| \sum_ {k = 1} ^ {n - 1} \lambda_ {k} a _ {k} \right\|.
\]

\begin{enumerate}
\def\labelenumi{\alph{enumi})}
\setcounter{enumi}{1}
\tightlist
\item
  设 \((a_{n})\) 是 E 的单位球中点的一个无穷序列，并设 \((\lambda_{n})\) 是数的一个无穷序列，诸数 \(>0\) 且满足 \(\sum_{n} \lambda_{n} = 1\) 。对每个 \(n > 0\) ，令 \(\mu_{n} = 1 - \sum_{k=1}^{n-1} \lambda_{k}\) ，并令 \(b_{n} = \sum_{k=1}^{n-1} \lambda_{k} a_{k} + \mu_{n} a_{n}\) ；证明：对所有 \(n \geqslant 1\) ，我们有
\end{enumerate}

\[
\left\| \sum_ {k = 1} ^ {n} \lambda_ {k} a _ {k} \right\| \leqslant \mu_ {n} \sum_ {k = 1} ^ {n} \frac {\lambda_ {k} \| b _ {k} \|}{\mu_ {k - 1} \mu_ {k}}.
\]

（用归纳法应用 a)。）

\begin{enumerate}
\def\labelenumi{\alph{enumi})}
\setcounter{enumi}{2}
\tightlist
\item
  设 \((\mathrm{C}_n)\) 是 E 中凸集的一个递减序列，包含于单位球中，且满足 \(d(0,\mathrm{C}_1)\geqslant \theta >0\) （从而更有 \(d(0,\mathrm{C}_n)\geqslant \theta\) 对所有 \(n\) ）。设 \((\lambda_{n})\) 是数的一个序列，诸数 \(>0\) 且满足 \(\sum_{n}\lambda_{n} = 1\) 。证明：存在一个数 \(\alpha\) 使得 \(\theta \leqslant \alpha \leqslant 1\) ，以及 E 的点的一个序列 \((x_{n})\) ，使得 \(x_{n}\in \mathbf{C}_{n}\) 对所有 \(n\) ，\(\| \sum_{n}\lambda_{n}x_{n}\| = \alpha\) ，且对所有 \(n\)
\end{enumerate}

\[
\left\| \sum_ {k = 1} ^ {n} \lambda_ {k} x _ {k} \right\| \leqslant \alpha (1 - \theta \sum_ {j = n + 1} ^ {\infty} \lambda_ {j}).
\]

（取 \(x_{1}\) 使 \(\| x_{1}\|\) 任意接近 \(d(0, \mathrm{C}_{1})\) ，然后用归纳法取 \(x_{n}\) ，使 \(\left\| \sum_{k=1}^{n-1} \lambda_{k} x_{k} + \left(\sum_{j=n}^{\infty} \lambda_{j}\right) x_{n}\right\|\) 任意接近诸数 \(\left\| \sum_{k=1}^{n-1} \lambda_{k} x_{k} + \left(\sum_{j=n}^{\infty} \lambda_{j}\right) y\right\|\) 的下确界，其中 \(y\) 取遍 \(\mathrm{C}_n\) 。然后利用 \(b)\) 。）

¶ 25) 设 E 是一个满足第一可数公理的 Banach 空间。证明下列性质等价：

α) E 不是自反的。

\(\beta)\) 对每个数 \(\theta\) （满足 \(0 < \theta < 1\) ），存在一个序列 \((x_n')\) （在 \(\mathbf{E}'\) 中），使得 \(\| x_n'\| \leqslant 1\) 对所有 \(n\) ，序列 \((x_n')\) 对于 \(\sigma (\mathrm{E}',\mathrm{E})\) 收敛于 0，并且 0 到由诸 \(x_{n}^{\prime}\) 生成的凸集的距离 \(\geqslant \theta\) 。

\(\gamma)\) 对每个数 \(\theta\) （满足 \(0 < \theta < 1\) ）以及每个序列 \(\lambda_{n}\) （诸数 \(>0\) ，满足 \(\sum_{n} \lambda_{n} = 1\) ），存在一个数 \(\alpha\) 使得 \(\theta \leqslant \alpha \leqslant 1\) ，以及点的一个序列 \((y_{n}')\) ，其点属于

E' ，使得 \(\| y_n'\| \leqslant 1\) 对所有 \(n\) ，\(\left\| \sum_{n}\lambda_ny_n'\right\| = \alpha\) ，且 \(\left\| \sum_{k=1}^{n}\lambda_ky_k'\right\| \leqslant \alpha (1 - \theta \sum_{j=n+1}^{\infty}\lambda_j)\) 对所有 \(n\) 。

\(\delta)\) 存在 \(z' \in \mathrm{E}'\) ，使得没有任何 \(x \in \mathrm{E}\) 满足 \(\left|\langle x, z' \rangle\right| = \| x \| \cdot \| z'\|\) （James-Klee 定理）。

（要看 \(\alpha\) ) 蕴含 \(\beta\) )，注意存在 \(z'' \in \mathrm{E}''\) 使得 \(\| z'' \| < 1\) 且 \(d(z'', \mathrm{E}) > \theta\) 。若 \((x_n)\) 是 \(\mathrm{E}\) 中的一个处处稠密序列，求一个序列 \((x_n')\) （在 \(\mathrm{E}'\) 中），使得 \(\| x_n' \| < 1\) ，且 \(\langle x_k, x_n' \rangle = 0\) （对 \(k \leqslant n\) ）以及 \(\langle x_n', z'' \rangle = \theta\) 。要看 \(\beta\) ) 蕴含 \(\gamma\) )，利用 exerc. 24。要看 \(\gamma\) ) 蕴含 \(\delta\) )，证明对每个 \(x \in \mathrm{E}\) 有 \(|\sum_{n} \lambda_n \langle x, y_n' \rangle| < \alpha\) （利用 \(\gamma\) ）。）

\begin{enumerate}
\def\labelenumi{\arabic{enumi})}
\setcounter{enumi}{25}
\tightlist
\item
  我们称局部凸空间 E 具有性质 (GDF)，如果从 E 到某个 Banach 空间 F 中的每个满足下列性质的线性映射 u 都是连续的：在乘积空间 \(E \times F\) 中，u 的图像 \(\Gamma\) 中收敛点列的每个极限仍属于 \(\Gamma\) 。每个 Fréchet 空间都具有性质 (GDF) (I, 第 19 页, 推论 5)；Fréchet 空间族的每个归纳极限也是如此 (II, 第 34 页, 命题 10)。证明：每个具有性质 (GDF) 的 Hausdorff 局部凸空间都是桶型的。（设 V 是 E 中的一个桶，q 是它的度规，H 是 E 带上这个半范数所对应的 Hausdorff 空间；证明典范映射 \(\pi\) （从 E 到完备化 \(\hat{H}\) 中的）是连续的，这要利用性质 (GDF) 以及每个线性型 \(x' \in V^{0}\) 都可以唯一地扩张成 \(\hat{H}\) 上的连续线性型这一事实（这些线性型所成的集是 \(\hat{H}\) 的对偶中的单位球）。）
\end{enumerate}

